\specialchapter{Bibliography}

\begin{enumerate}
\def\labelenumi{\arabic{enumi}.}
\item
  L. EULER, Vollständige Anleitung zur Algebra (=Opera Omnia (1), vol.~I, Leipzig-Berlin (Teubner), 1911).
\item
  J. L. LAGRANGE, Oeuvres, 14 volumes, Paris (Gauthier-Villars), 1867--1892.
\item
  C. F. GAUSS, Werke, 12 volumes, Göttingen, 1870--1927.
\item
  P. G. LEJEUNE-DIRICHLET, Werke, 2 volumes, Berlin (Reimer), 1889-1897.
\end{enumerate}

4 (bis). P. G. LEJEUNE-DIRICHLET, Vorlesungen über Zahlentheorie, 2te Aufl., Braunschweig (Vieweg), 1871.

\begin{enumerate}
\def\labelenumi{\arabic{enumi}.}
\setcounter{enumi}{4}
\item
  C. G. J. JACOBI, Gesammelte Werke, 7 volumes, Berlin (Reimer), 1881-1891.
\item
  G. EISENSTEIN: (a) Beweis der Reciprocitätsgesetze für die cubischen Reste in der Theorie der aus dritten Wurzeln der Einheit zusammengesetzen Zahlen, Crelle's Journal, 27 (1844), pp.~289--310; (b) Zur Theorie der quadratischen Zerfällung der Primzahlen 8n + 3, 7n + 2 und 7n + 4, Crelle's Journal, 37 (1848), pp.~97--126; (c) Über einige allgemeine Eigenschaften der Gleichung von welcher die Teilung der ganzen Lemniscate abhängt, nebst Anwendungen derselben auf die Zahlentheorie, Crelle's Journal, 39 (1850), pp.~160--179 and 224--287.
\item
  E. KUMMER: (a) Sur les nombres complexes qui sont formés avec les nombres entiers réels et les racines de l'unité, J. de Math., (1), 12 (1847), pp.~185--212; (b) Zur Theorie der complexen Zahlen, Crelle's Journal, 35 (1847), pp.~319--326; (c) Ueber die Zerlegung der aus Wurzeln der Einheit gebildeten complexen Zahlen in Primfactoren, Crelle's Journal, 35 (1847), pp.~327--367; (d) Mémoire sur les nombres complexes composés de racines de l'unité et des nombres entiers, J. de Math., (1), 16 (1851), pp.~377--498; (e) Über die allgemeinen Reciprocitätsgesetze unter den Resten und Nichtresten der Potenzen deren Grad eine Primzahl ist (Abh. der Kön. Akad. der Wiss. zu Berlin (1859), Math. Abhandl., pp.~19--159).
\item
  C. HERMITE, Oeuvres, 4 Volumes, Paris (Gauthier-Villars), 1905--1917.
\item
  L. KRONECKER, Werke, 5 volumes, Leipzig (Teubner), 1895--1930: (a) De unitatibus complexis, vol.~I, pp.~5--71 (=Inaug. Diss., Berolini, 1845); (b) Über die algebraisch auflösbaren Gleichungen I, vol.~IV, pp.~1--11 (=Monatsber. der Kön. Preuss. Akad. der Wiss., 1853, pp.~365--374); (c) Über die elliptischen Functionen für welche complexe Multiplication stattfindet vol.~IV, pp.~177--183 (=Monatsber. der Kön. Preuss. Akad. der Wiss., 1857, pp.~455--460); (d) Über die complexe Multiplication der elliptischen Functionen, vol.~IV, pp.~207--217 (=Monatsber. der Kön. Preuss. Akad. der Wiss., 1862, pp.~363--372); (e) Über die Discriminante algebraischer Functionen einer Variabeln, vol.~II, pp.~193--236 (=Crelle's Journal, 91 (1881), pp.~301--334); (f) Grundzüge einer arithmetischen Theorie der
\end{enumerate}

algebraischen Grössen, vol.~II, pp.~237--387 (=Crelle's Journal, 92 (1882), pp.~1--122).

\begin{enumerate}
\def\labelenumi{\arabic{enumi}.}
\setcounter{enumi}{9}
\tightlist
\item
  R. DEDEKIND, Gesammelte mathematische Werke, 3 volumes, Braunschweig (Vieweg), 1932: (a) Abriss einer Theorie der höheren Kongruenzen in bezug auf einen reellen Primzahl-Modulus, vol.~I, pp.~40--66 (=Crelle's Journal, 54 (1857), pp.~1--26; (b) Sur la Théorie des Nombres entiers algébriques, vol.~III, pp.~262--296 (=Bull. Sci. Math., (1), 11 (1876), pp.~278--288 and (2), 1 (1877), pp.~17--41, 69--92, 144--164, 207--248); (c) Über die Anzahl der Ideal-Klassen in den verschiedenen Ordnungen eines endlichen Körpers, vol.~I, pp.~105--157 (=Festschrift der Technischen Hochschule in Braunschweig zur Säkularfeier des Geburtstages von C. F. Gauss, Braunschweig, 1877, pp.~1--55); (d) Über den Zusammenhang zwischen der Theorie der Ideals und der Theorie der höheren Kongruenzen, vol.~I, pp.~202--230 (=Abh. Kön. Ges. Wiss. zu Göttingen, 23 (1878), pp.~1--23); (e) Über die Discriminanten endlicher Körper, vol.~I, pp.~351--396 (=Abh. Kön. Ges. Wiss. zu Göttingen, 29 (1882), pp.~1--56); (f) Über die Theorie der ganzen algebraischen Zahlen, vol.~III, pp.~1--222 (=Supplement XI von Dirichlets Vorlesungen über Zahlentheorie, 4 Aufl. (1894), pp.~434--657); (g) Zur Theorie der Ideale, vol.~II, pp.~43--48 (=Nachr. Göttingen, 1894, pp.~272--277); (h) Über eine Erweiterung des Symbols (a, b) in der Theorie der Moduln, vol.~II, pp.~59--85 (=Nachr. Göttingen, 1895, pp.~183--208).
\end{enumerate}

10 (bis). R. DEDEKIND-H. WEBER, Theorie der algebraischen Funktionen einer Veränderlichen, Crelle's Journal, 92 (1882), pp.~181--290 (=R. Dedekind, Ges. Math. Werke, vol.~I, pp.~238--349).

\begin{enumerate}
\def\labelenumi{\arabic{enumi}.}
\setcounter{enumi}{10}
\item
  E. SELLING, Ueber die idealen Primfactoren der complexen Zahlen, welche aus den Wurzeln einer beliebigen irreductiblen Gleichung rational gebildet sind, Zeitschr. für Math. und Phys., 10 (1865), pp.~17--47.
\item
  M. NOETHER, Über einen Satz aus der Theorie der algebraischen Funktionen, Math. Ann., 6 (1873), pp.~351--359.
\item
  A. BRILL-M. NOETHER, Ueber algebraischen Funktionen, Math. Ann., 7 (1874), pp.~269--310.
\item
  G. ZOLOTAREFF, Sur la théorie des nombres complexes, J. de Math. (3), 6 (1880), pp.~51--84 and 129--166.
\item
  E. NETTO, Zur Theorie der Elimination, Acta Math., 7 (1885), pp.~101-104.
\item
  D. HILBERT: (a) Über die Theorie der algebraischen Formen, Math. Ann., 36 (1890), pp.~473--534; (b) Über die vollen Invariantensysteme, Math. Ann., 42 (1893), pp.~313--373; (c) Grundzüge einer Theorie des Galoischen Zahlkörpers, Gött. Nachr., (1894), pp.~224--236; (d) Zahlbericht, Jahresber. der D. M. V., 4 (1897), pp.~175--546 (translated into French by A. Lévy and Th. Got under the title ``Théorie des corps de nombres algébriques'', Paris (Hermann), 1913).
\item
  K. WEIERSTRASS, Mathematische Werke, 7 volumes, Berlin (Mayer und Müller), 1894--1927.
\item
  K. HENSEL: (a) Über eine neue Begründung der Theorie der algebraischen Zahlen, Jahresber. der D. M. V., 6 (1899), pp.~83--88; (b) Ueber die Fundamentalgleichung und die ausserwesentlichen Diskriminantentheiler eines algebraischen Körpers, Gött. Nachr., (1897), pp.~254--260; (c) Neue Grundlagen der Arithmetik, Crelle's Journal, 127 (1902), pp.~51--84; (d) Über die arithmetische Eigenschaften der algebraischen und transzendenten Zahlen, Jahresber. der D. M. V., 14 (1905), pp.~545--558; (e) Ueber die arithmetischen Eigenschaften der Zahlen, Jahresber. der D. M. V., 16 (1907), pp.~299--319, 388--393, 474--496; (f) Theorie der algebraischen Zahlen, Leipzig (Teubner), 1908.
\item
  E. LASKER, Zur Theorie der Moduln und Ideale, Math. Ann., 60 (1905), pp.~20--116.
\item
  E. STEINITZ: (a) Algebraische Theorie der Körper, Crelle's Journal, 137 (1910), pp.~167--308; (b) Rechteckige Systeme und Moduln in algebraischen Zahlkörpern, Math. Ann. 71 (1912), pp.~328--354 and 72 (1912), pp.~297--345.
\item
  F. S. MACAULAY, On the resolution of a given modular system into primary systems including some properties of Hilbert numbers, Math. Ann., 74 (1913), pp.~66--121.
\item
  J. KÜRSCHAK, Über Limesbildung und allgemeine Körpertheorie, Crelle's Journal, 142 (1913), pp.~211--253.
\item
  A. FRAENKEL, Über die Teiler der Null und die Zerlegung von Ringen, Crelle's Journal, 145 (1914), pp.~139--176.
\item
  A. OSTROWSKI, Über einige Lösungen der Funktionalgleichung \(\phi(x)\phi(y) = \phi(x.y)\) , Acta Math., 41 (1917), pp.~271-284.
\item
  E. NOETHER: (a) Idealtheorie in Ringbereichen, Math. Ann., 83 (1921), pp.~24--66; (b) Abstrakter Aufbau der Idealtheorie in algebraischen Zahl- und Funktionenkörpern, Math. Ann., 96 (1927), pp.~26--61.
\item
  H. HASSE: (a) Ueber die Darstellbarkeit von Zahlen durch quadratischen Formen im Körper der rationalen Zahlen, Crelle's Journal, 152 (1923), pp.~129--148; (b) Ueber die Äquivalenz quadratischer Formen im Körper der rationalen Zahlen, Crelle's Journal, 152 (1923), pp.~205--224; (c) Kurt Hensels entscheidender Anstoss zur Entdeckung des Lokal-Global-Prinzips, Crelle's Journal, 209 (1960), pp.~3--4.
\item
  H. JUNG, Algebraischen Flächen, Hannover (Helwing), 1925.
\item
  H. GRELL, Beziehung zwischen den Idealen verschiedener Ringe, Math. Ann., 97 (1927), pp.~490-523.
\item
  W. KRULL: (a) Primidealketten in allgemeine Ringbereichen, Sitz. Ber. Heidelberg Akad. Wiss., 1928; (b) Allgemeine Bewertungstheorie, Crelle's Journal, 167 (1931), pp.~160--196; (c) Beiträge zur Arithmetik kommutativer Integritätsbereiche, III, Math. Zeitschr., 42 (1937), pp.~745--766; (d) Dimensionstheorie in Stellenringen, Crelle's Journal, 179 (1938), pp.~204--226; (e) Idealtheorie, Berlin (Springer), 1935.
\item
  M. H. STONE: (a) The theory of representations for Boolean algebras, Trans. Amer. Math. Soc., 40 (1936), pp.~37-111; (b) Applications of the theory of Boolean rings to general topology, Trans. Amer. Math. Soc., 41 (1937), pp.~375-481.
\item
  B. L. van der WAERDEN, Moderne Algebra, vol.~II, Berlin (Springer), 1931.
\item
  C. CHEVALLEY: (a) Généralisation de la théorie du corps de classes pour les extensions infinies, J. de Math., (9), 15 (1936), pp.~359--371; (b) On the theory of local rings, Ann. of Math., 44 (1943), pp.~690--708.
\item
  O. ZARISKI: (a) The compactness of the Riemann manifold of an abstract field of algebraic functions, Bull. Amer. Math. Soc., 50 (1944), pp.~683--691; (b) Generalized semi-local rings, Summa Bras. Math., 1 (1946), pp.~169--195.
\item
  N. JACOBSON, A topology for the set of primitive ideals in an arbitrary ring, Proc. Nat. Acad. Sci. U.S.A., 31 (1945), pp.~333--338.
\item
  I. COHEN-A.SEIDENBERG, Prime ideals and integral independence, Bull. Amer. Math. Soc., 52 (1946), pp.~252--261.
\item
  P. SAMUEL, La notion de multiplicité en Algèbre et en Géométrie algébrique, J. de Math., (9), 30 (1951), pp.~159--274.
\item
  A. WEIL, Fibre-spaces in Algebraic Geometry (Notes by A. Wallace), Chicago Univ., 1952.
\item
  J. P. SERRE: (a) Faisceaux algébriques cohérents, Ann. of Math., 61 (1955), pp.~197--278; (b) Géométrie algébrique et géométrie analytique, Ann. Inst. Fourier, 6 (1956), pp.~1--42.
\item
  A. GROTHENDIECK, Éléments de géométrie algébrique, Publ. math. Inst. Htes Et. Scient., 1960.
\end{enumerate}

